\section{Code, repository and datasets}
\label{sec:appendix}

\noindent\textbf{Our GitHub repository:} \url{\GitHubURL}

\medskip
\noindent\textbf{Datasets used.}
\begin{itemize}
\item Our own recorded dataset (15,378 labelled hand-point frames and 7,692
hand-crop images from 3 users, 7 gestures), in the repository under
\texttt{data/}.
\item Related public dataset (not used for training, listed for reference):
HaGRID~\cite{kapitanov2024hagrid}, \url{https://github.com/hukenovs/hagrid}.
\end{itemize}

\noindent\textbf{Repositories and libraries used.}
\begin{itemize}
\item MediaPipe (hand landmarks): \url{https://github.com/google-ai-edge/mediapipe}
\item OpenCV (camera and drawing): \url{https://github.com/opencv/opencv}
\item PyTorch (neural networks): \url{https://github.com/pytorch/pytorch}
\item scikit-learn (SVM, Random Forest, kNN baselines): \url{https://github.com/scikit-learn/scikit-learn}
\item PyAutoGUI (mouse and keyboard events): \url{https://github.com/asweigart/pyautogui}
\item python-pptx (slides): \url{https://github.com/scanny/python-pptx}
\item CVPR author kit (this template): \url{https://github.com/cvpr-org/author-kit}
\item One-Euro filter reference page: \url{https://gery.casiez.net/1euro/}
\end{itemize}

\noindent\textbf{How to run.} \texttt{pip install -r requirements.txt};
record data with \texttt{python collect\_data.py --user NAME}; run every
experiment with \texttt{python run\_all.py}; start the app with
\texttt{python app.py}. Details are in \texttt{README.md}.

\section{Complete source code}
\subsection*{\texttt{app.py}}
\lstinputlisting{code/app.py}
\subsection*{\texttt{collect\_data.py}}
\lstinputlisting{code/collect_data.py}
\subsection*{\texttt{config.py}}
\lstinputlisting{code/config.py}
\subsection*{\texttt{controller.py}}
\lstinputlisting{code/controller.py}
\subsection*{\texttt{features.py}}
\lstinputlisting{code/features.py}
\subsection*{\texttt{filters.py}}
\lstinputlisting{code/filters.py}
\subsection*{\texttt{hand\_tracker.py}}
\lstinputlisting{code/hand_tracker.py}
\subsection*{\texttt{make\_figures.py}}
\lstinputlisting{code/make_figures.py}
\subsection*{\texttt{make\_report.py}}
\lstinputlisting{code/make_report.py}
\subsection*{\texttt{make\_slides.py}}
\lstinputlisting{code/make_slides.py}
\subsection*{\texttt{make\_slides\_full.py}}
\lstinputlisting{code/make_slides_full.py}
\subsection*{\texttt{model.py}}
\lstinputlisting{code/model.py}
\subsection*{\texttt{personalize.py}}
\lstinputlisting{code/personalize.py}
\subsection*{\texttt{rules.py}}
\lstinputlisting{code/rules.py}
\subsection*{\texttt{run\_all.py}}
\lstinputlisting{code/run_all.py}
\subsection*{\texttt{train.py}}
\lstinputlisting{code/train.py}
\subsection*{\texttt{train\_cnn.py}}
\lstinputlisting{code/train_cnn.py}
\subsection*{\texttt{virtual\_keyboard.py}}
\lstinputlisting{code/virtual_keyboard.py}
\subsection*{\texttt{voting.py}}
\lstinputlisting{code/voting.py}
\subsection*{\texttt{experiments/analyze\_logs.py}}
\lstinputlisting{code/experiments__analyze_logs.py}
\subsection*{\texttt{experiments/click\_test.py}}
\lstinputlisting{code/experiments__click_test.py}
\subsection*{\texttt{experiments/typing\_test.py}}
\lstinputlisting{code/experiments__typing_test.py}

